% 2.	Formulação do Problema de Controle de Órbita e Atitude para uma  Espaçonave tipo Manipulador Robótico em Missão de RVD/B

% 2.1.	Formulação do Problema de Controle
% 2.1.1.	Definição de Estados e Variáveis de Controle
% 2.1.2.	Restrições Operacionais
% 2.1.3.	Critérios de Desempenho

\section{Formulação do Problema de Controle de Órbita e Atitude para uma Espaçonave tipo Manipulador Robótico em Missão de RVD/B}


\subsection{Formulação do Problema de Controle}


\subsubsection{Definição de Estados e Variáveis de Controle}


\subsubsection{Restrições Operacionais}

Tabela com os centros de massa dos elementos do manipulador robótico, conforme a tabela \ref{tab:centrosMassa_componentes}.
\begin{table}[h]
\centering
\caption{Centros de massa (em metros)}
\label{tab:centrosMassa_componentes}
\begin{tabular}{lrrr}
\hline
\textbf{Elemento} & \textbf{$x$ [m]} & \textbf{$y$ [m]} & \textbf{$z$ [m]}\\
\hline
Base geral                 &  0.000000 &  0.000000 &  0.050000 \\
Disco                      &  0.000000 & -0.225000 &  0.105000 \\
Junta 1                    &  0.000000 & -0.225000 &  0.112500 \\
Elo 1                      &  0.000000 & -0.225000 &  0.165000 \\
Junta 2                    &  0.000000 & -0.197500 &  0.165000 \\
Elo 2                      &  0.000000 & -0.170000 &  0.315000 \\
Junta 3                    &  0.000000 & -0.142500 &  0.465000 \\
Elo 3                      &  0.000000 & -0.115000 &  0.615000 \\
Junta 4                    &  0.000000 & -0.087500 &  0.765000 \\
Elo 4                      & -0.000010 & -0.059960 &  0.915030 \\
Junta 5                    & -0.000010 & -0.059960 &  1.067530 \\
Ferramenta                 & -0.000010 & -0.059960 &  1.145030 \\
Manipulador (global)       &  0.000000 & -0.004950 &  0.066530 \\
\hline
\end{tabular}
\end{table}

Tabela com os momentos de inércia dos elementos do manipulador robótico, conforme a tabela \ref{tab:matrizInercia_componentes}, todos em kg·m².

\begin{table}[h]
\centering
\footnotesize
\caption{Componentes da matriz de inércia \(\mathbf{I}\) de cada subsistema (kg·m$^2$).}
\label{tab:matrizInercia_componentes}
\begin{tabular}{lrrrrrr}
\toprule
Componente & $I_{xx}$ & $I_{yy}$ & $I_{zz}$ & $I_{xy}$ & $I_{xz}$ & $I_{yz}$ \\
\midrule
Base Geral              & 3.274000 & 3.274000 & 6.441000 & 0.000000 & 0.000000 & 0.000000 \\
Disco                   & 0.000788 & 0.000788 & 0.001571 & 0.000000 & 0.000000 & 0.000000 \\
Junta 1                 & 0.000002 & 0.000002 & 0.000003 & 0.000000 & 0.000000 & 0.000000 \\
Elo 1                   & 0.000194 & 0.000194 & 0.000061 & 0.000000 & 0.000000 & 0.000000 \\
Junta 2                 & 0.000002 & 0.000003 & 0.000002 & 0.000000 & 0.000000 & 0.000000 \\
Elo 2                   & 0.004510 & 0.004510 & 0.000184 & 0.000000 & 0.000000 & 0.000000 \\
Junta 3                 & 0.000002 & 0.000003 & 0.000002 & 0.000000 & 0.000000 & 0.000000 \\
Elo 3                   & 0.004910 & 0.004510 & 0.000184 & 0.000000 & 0.000000 & 0.000000 \\
Junta 4                 & 0.000002 & 0.000003 & 0.000002 & 0.000000 & 0.000000 & 0.000000 \\
Elo 4                   & 0.004510 & 0.004510 & 0.000184 & 0.000000 & 0.000000 & 0.000000 \\
Junta 5                 & 0.000002 & 0.000002 & 0.000003 & 0.000000 & 0.000000 & 0.000000 \\
Ferramenta              & 0.000050 & 0.000050 & 0.000001 & 0.000000 & 0.000000 & 0.000000 \\
Manipulador (global)    & 4.045000 & 3.993000 & 6.496000 & 0.000000 & 0.000003 & 0.102800 \\
\bottomrule
\end{tabular}

\vspace{0.3em}
\footnotesize{Observação: produtos de inércia são nulos em todos os subsistemas exceto no MI global, onde \(I_{xz}\) e \(I_{yz}\) são não nulos.}
\end{table}

\newpage

Tabela com os parâmetros físicos dos elementos do manipulador robótico, conforme a tabela \ref{tab:parametros_fisicos}, todos em kg.

\begin{table}[h]
\centering
\caption{Parâmetros físicos dos subsistemas do manipulador.}
\label{tab:parametros_fisicos}
\begin{tabular}{lrrr}  % ← Corrigido para 4 colunas
\toprule
Componente & Altura [m] & Diâmetro [m] & Massa [kg] \\
\midrule
Base Geral  & 0.100 & 0.900 & 63.6173 \\
Disco       & 0.010 & 0.200 & 0.3142 \\
Junta 1     & 0.005 & 0.050 & 0.0098 \\
Elo 1       & 0.100 & 0.050 & 0.1964 \\
Junta 2     & 0.005 & 0.050 & 0.0098 \\
Elo 2       & 0.300 & 0.050 & 0.5890 \\
Junta 3     & 0.005 & 0.050 & 0.0098 \\
Elo 3       & 0.300 & 0.050 & 0.5890 \\
Junta 4     & 0.005 & 0.050 & 0.0098 \\
Elo 4       & 0.300 & 0.050 & 0.5890 \\
Junta 5     & 0.005 & 0.050 & 0.0098 \\
Ferramenta  & 0.150 & 0.015 & 0.0265 \\
\midrule
\textbf{Total (manipulador)} & --- & --- & \textbf{65.9704} \\
\bottomrule
\end{tabular}
\end{table}







\subsubsection{Critérios de Desempenho}


